% ============================================================
% 06_conclusion.tex — Conclusion (Macedonian)
% ============================================================

\section{Резиме на архитектурата и придонесите}

Оваа дипломска работа го презентираше проектирањето, имплементацијата и евалуацијата на целосно автоматизиран систем за управување со FPL тим базиран на тројно-слојна хибридна архитектура на вештачка интелигенција. Секој слој на архитектурата дава специфичен придонес:

\textbf{Слој 1 — LightGBM предикција}: Четирите позициски модели обезбедуваат квантитативни предикции на поените базирани на богат сет од историски и тековни карактеристики. Walk-forward валидацијата обезбеди дека моделите се валидирани на неостварен временски период, гарантирајќи реалистични проценки на точноста.

\textbf{Слој 2 — ILP оптимизација}: Формулацијата на проблемот за избор на тим како ILP дава математичка гаранција за оптималноста на изборот во рамките на дефинираните ограничувања. Управувањето со чипови и трансферната стратегија се интегрирани директно во оптимизаторот, обезбедувајќи кохерентни долгорочни одлуки.

\textbf{Слој 3 — Intel Pipeline и LLM}: Седумте intel модули обезбедуваат pre-deadline разузнавачки сигнали кои ги коригираат предикциите врз основа на информации кои не можат да се пронајдат во историскиот датасет: тековни повреди, тренерски изјави, пазарни движења. Нарацискиот слој преку Claude API ги прави одлуките разбирливи за крајниот корисник.

\section{Клучни наоди}

Евалуацијата на системот открила неколку значајни наоди:

\textbf{LightGBM доминира над XGBoost}: Во рамките на joint random search со 250 обиди, 14 од 15-те најдобри конфигурации користеле LightGBM. Разликата во перформансите е конзистентна низ сите позиции и набори, потврдувајќи дека leaf-wise стратегијата за раст на дрвата е поефикасна за оваа задача.

\textbf{Intel pipeline додава значителна вредност}: Додавањето на intel казни за достапноста и ризикот за ротација придонесе за подобрување од $\approx$59 поени на GW1-10 размер. Ова е особено важно бидејќи информациите за повреди и ротации не можат да се изведат само од историскиот датасет.

\textbf{Bayesian оптимизацијата е поефикасна}: Optuna TPE постигна +39 поени над случајното пребарување со само 40\% од бројот на обиди, потврдувајќи ја практичноста на Bayesian методите за оптимизација на хиперпараметри при скапи функции на цел.

\textbf{Откривање и поправање на грешки е круцијално}: Пет критични грешки беа идентификувани и поправени во текот на развојот. Без систематско тестирање и евалуација, овие грешки би останале незабележани и би предизвикале погрешни конклузии за перформансите на системот.

\section{Финален резултат и значење}

Системот постигна \textbf{1799 поени} во текот на GW1--28 (просек \textbf{64.3 поени/ГН}) со нула казнени поени. Ова е постигнато во строги услови на blind test (нула информации за тековната сезона при тренирање), со потпирање исклучиво на историски податоци и real-time intel сигнали пред секој рок.

Претставениот систем е, доколку се знае за авторите, прв систем кој ги комбинира сите следни компоненти во единствен FPL управувачки пипелајн: LightGBM ML предикција со позициска специјализација, ILP оптимизација со полна поддршка за чипови, pre-deadline intel pipeline со 7 модули, и LLM нарациски слој. Ова претставува значаен чекор напред во однос на постоечките пристапи кои обично се фокусираат само на еден или два аспекта.

\section{Завршни мисли}

Управувањето со FPL тим е мулти-фазен, мулти-ограничувачки проблем на одлучување под неизвесност — точно видот на проблем за кој хибридните AI системи се особено добро позиционирани. Оваа теза покажала дека со соодветна комбинација на ML, оптимизација и интелигентно собирање на информации, е можно да се надмине просечниот FPL менаџер и да се донесуваат рационални, объаснети одлуки за целата сезона.

Работата останува отворена за натамошни подобрувања — нови карактеристики, live деплојмент, RL пристапи — но основниот систем е функционален, евалуиран, и готов за продукциска употреба.
